%-------------------------------------------------------------------------------
%	SECTION TITLE
%-------------------------------------------------------------------------------
\cvsection{Projects}


%-------------------------------------------------------------------------------
%	CONTENT
%-------------------------------------------------------------------------------
\begin{cventries}
	
	%---------------------------------------------------------
	\cventry
	{vLLM, TGI, Ollama, Kubernetes, LangChain, LangGraph, Chroma} % Tech stack
	{Local GenAI Stack - LLM Serving, RAG, and Agentic Orchestration} % Title
	{}
	{}
	{
		\begin{cvitems}
			\item {Deployed Qwen, DeepSeek and LLaMA-family open-weight models fully offline on a self-managed Kubernetes cluster, benchmarking vLLM, Hugging Face TGI and Ollama as inference runtimes.}
		    \item {Exposed an OpenAI-compatible inference API through vLLM and configured tool calling with automatic tool choice and a Hermes-style parser for agentic applications.}
		   \item {Built LangGraph orchestration workflows with stateful multi-step execution, tool selection, conditional routing and handoff patterns.}
	       	\item {Ran GGUF, GPTQ and AWQ quantized models and measured throughput/latency trade-offs across runtimes and quantization levels to select suitable deployment configurations for constrained GPU resources.}
	    	\item {Built a LangChain RAG pipeline over a Markdown knowledge base using local embeddings and Chroma, tuning chunking, retrieval top-K and prompt construction against context-window constraints.}
	    \end{cvitems}
	}

	%---------------------------------------------------------
    
    \cventry
    {Kubernetes, Proxmox, Helm, Prometheus, Grafana, Loki, Tailscale} % Tech stack
    {Self-Hosted Kubernetes Platform } % Title
    {}
    {}
    {
    	\begin{cvitems}
    		\item Built and operate a 3-node kubeadm Kubernetes cluster on Proxmox and bare-metal nodes, running services such as file sync, document management, self-hosted AI model stack and Forgejo, a self-hosted Git.
    		\item Designed the storage layer with dynamic PV provisioning through CSIs, CloudNativePG for stateful Postgres workloads, and backups.
    		\item Implemented cluster networking and access control, along with Tailscale to manage remote connectivity to the cluster applications
    		\item Deployed full-stack observability with kube-prometheus-stack, Grafana, Loki and Alloy for metrics and logs across cluster and application workloads.
    	\end{cvitems}
    }
	
	%---------------------------------------------------------
	\cventry
	{PyTorch, TensorFlow, Scikit-learn, NLP} % Tech stack
	{Advertisement Detection (Final-Year Capstone)} % Title
	{}
	{}
	{
		\begin{cvitems}
			\item Personally implemented and trained classifiers on a 1M+ row URL dataset assembled from multiple sources
			\item Engineered NLP-based features from URL structure to detect ad links.
			\item Compared neural networks, XGBoost, SVMs, LightGBM and Random Forests to find the best-performing model and reached 80\% classification accuracy.
		\end{cvitems}
	}
	
	%---------------------------------------------------------
	%	\cventry
	%	{TensorFlow, Scikit-learn, NLP, Pandas} % Tech stack
	%	{Sentiment Analysis - Amazon Reviews (227k) and Twitter (10k+)} % Title
	%	{}
	%	{}
	%	{
	%		\begin{cvitems}
	%			\item Two end-to-end NLP pipelines: ETL, text preprocessing, feature engineering, model training, evaluation and visualisation.
	%			\item Achieved 90\% classification accuracy on Amazon and 92\% on Twitter by combining neural networks with classical methods (XGBoost, SVMs, Random Forests) and hyperparameter tuning.
	%		\end{cvitems}
	%	}
	
	
	%---------------------------------------------------------
	\cventry
	{Python, REST APIs, PyPi, Github CI/CD pipelines} % Tech stack
	{vscan - Open-source Python security scanner} % Title
	{}
	{}
	{
		\begin{cvitems}
			\item Project lead on a Python-based virus scanner that uses different APIs to scan files and websites for malicious content.
			\item Built a GitHub Actions CI/CD pipeline that runs tests on every pull request and automatically builds and publishes versioned releases to PyPi on tagged commits.
		\end{cvitems}
	}
	
	%---------------------------------------------------------
\end{cventries}
